\section{Algoritmus RSA}\label{sec:krypto-rsa}

RSA je nejznámější příklad asymetrického kryptografického systému. Jeho název vznikl z~příjmení autorů Rona \textbf{R}ivesta, Adiho \textbf{S}hamira a~Leonarda \textbf{A}dlemana. Základní operací je umocňování modulo součin dvou velkých prvočísel.

\subsection{Vytvoření klíčů}

Bob postupuje následovně:
\begin{enumerate}
    \item zvolí dvě různá velká prvočísla $p$ a~$q$ a~spočítá
    \[
        n=pq;
    \]
    \item spočítá
    \[
        \varphi(n)=(p-1)(q-1);
    \]
    \item zvolí číslo $s$ nesoudělné s~$\varphi(n)$;
    \item nalezne číslo $v$ takové, že
    \[
        vs\bmod\varphi(n)=1;
    \]
    \item zveřejní klíč $k_V=(v,n)$ a~uchová v~tajnosti klíč $k_S=(s,n)$.
\end{enumerate}

Číslo $v$ je multiplikativní inverze čísla $s$ modulo $\varphi(n)$. Existuje právě proto, že $\operatorname{gcd}(s,\varphi(n))=1$, a~lze je efektivně nalézt rozšířeným Eukleidovým algoritmem.

\warningbox{Číslo $s$ je soukromý exponent, nikoli třetí prvočíslo. Musí být nesoudělné s~$(p-1)(q-1)$, ale samo prvočíslem být nemusí.}

\subsection{Šifrování a~dešifrování}

Zprávu reprezentujeme číslem $x$ splňujícím $0\leqslant x<n$. Alice zašifruje
\[
    y=\mathbf{E}(x,k_V)=x^v\bmod n.
\]
Bob dešifruje
\[
    \mathbf{D}(y,k_S)=y^s\bmod n.
\]
Mocniny mohou být obrovské, v~programu je však nepočítáme nejprve jako běžná celá čísla. Opakovaným umocňováním a~průběžným výpočtem zbytku lze modulární mocninu určit v~čase polynomiálním vzhledem k~počtu bitů vstupu.

\begin{figure}[ht]
    \centering
    \input{06-kryptografie/images/ch06-rsa-schema.tex}
    \caption{Schéma šifrování a~dešifrování v~RSA.}
    \label{fig:krypto-rsa-schema}
\end{figure}

\subsection{Malá Fermatova věta a~čínská věta o~zbytcích}

Korektnost RSA odvodíme ze dvou vět elementární teorie čísel.

\begin{theorem}[Malá Fermatova věta]\label{thm:krypto-fermat}
    Nechť $p$ je prvočíslo a~$a\in\Z$ není dělitelné číslem $p$. Potom
    \[
        a^{p-1}\bmod p=1.
    \]
    Ekvivalentně pro každé $a\in\Z$ platí $a^p\bmod p=a\bmod p$.
\end{theorem}

\begin{theorem}[Čínská věta o~zbytcích]\label{thm:krypto-cinska}
    Nechť $m_1,\ldots,m_k\in\N$ jsou po dvou nesoudělná. Pro libovolná $a_1,\ldots,a_k\in\Z$ existuje řešení soustavy
    \[
        x\equiv a_i\pmod{m_i},
        \qquad i=1,\ldots,k,
    \]
    a~všechna její řešení jsou navzájem shodná modulo $M=m_1\cdots m_k$.
\end{theorem}

\begin{corollary}\label{cor:krypto-cinska-dve}
    Jsou-li $p$ a~$q$ různá prvočísla a~platí
    \[
        a\equiv b\pmod p
        \qquad\text{a}\qquad
        a\equiv b\pmod q,
    \]
    potom $a\equiv b\pmod{pq}$.
\end{corollary}

\subsection{Korektnost RSA}

\begin{theorem}\label{thm:krypto-rsa-korektnost}
    Pro klíče vytvořené uvedeným postupem a~každé celé číslo $x$ splňující $0\leqslant x<n$ platí
    \[
        \mathbf{D}\bigl(\mathbf{E}(x,k_V),k_S\bigr)=x.
    \]
\end{theorem}

\begin{proof}
    Z~rovnosti $vs\bmod(p-1)(q-1)=1$ plyne, že pro nějaké $t\in\N_0$ platí
    \[
        vs=1+t(p-1)(q-1).
    \]
    Nejprve ukážeme $x^{vs}\equiv x\pmod p$. Pokud $p$ dělí $x$, jsou obě strany kongruence nulové modulo $p$. Pokud $p$ číslo $x$ nedělí, použijeme \smartref[větu~][]{\showrefs}{thm:krypto-fermat}{malou Fermatovu větu}:
    \begin{align*}
        x^{vs}
        &=
        x^{1+t(p-1)(q-1)}\\
        &=
        x\left(x^{p-1}\right)^{t(q-1)}
        \equiv
        x\cdot 1^{t(q-1)}
        \equiv x
        \pmod p.
    \end{align*}
    Stejným rozlišením případů dostaneme
    \[
        x^{vs}\equiv x\pmod q.
    \]
    Podle \smartref[důsledku~][]{\showrefs}{cor:krypto-cinska-dve}{důsledku čínské věty o~zbytcích} tedy
    \[
        x^{vs}\equiv x\pmod{pq}.
    \]
    Protože $n=pq$ a~$0\leqslant x<n$, uzavíráme
    \[
        \mathbf{D}\bigl(\mathbf{E}(x,k_V),k_S\bigr)
        =
        \left(x^v\bmod n\right)^s\bmod n
        =
        x^{vs}\bmod n
        =
        x.
    \]
\end{proof}

\begin{figure}[ht]
    \centering
    \begin{tikzpicture}[
    font=\small,
    >=Latex,
    box/.style={rounded corners=1.2mm, draw=black!55, fill=black!4,
                minimum width=3.3cm, minimum height=9mm, align=center},
    residue/.style={rounded corners=1.2mm, draw=blue!65!black, fill=blue!11,
                    minimum width=4.1cm, minimum height=10mm, align=center},
    result/.style={rounded corners=1.2mm, draw=green!55!black, fill=green!12,
                   minimum width=5cm, minimum height=10mm, align=center},
    flow/.style={-{Latex}, draw=black!70, line width=.9pt}
]
    \node[box] (power) at (0,2.2) {\(x^{vs}\), kde \(vs=1+t(p-1)(q-1)\)};
    \node[residue] (p) at (-3.3,.25)
        {\(x^{vs}\equiv x\pmod p\)\\malá Fermatova věta};
    \node[residue] (q) at (3.3,.25)
        {\(x^{vs}\equiv x\pmod q\)\\malá Fermatova věta};
    \node[result] (pq) at (0,-2)
        {\(x^{vs}\equiv x\pmod{pq}\)\\čínská věta o~zbytcích};

    \draw[flow] (power.south west) -- (p.north);
    \draw[flow] (power.south east) -- (q.north);
    \draw[flow] (p.south) -- (pq.north west);
    \draw[flow] (q.south) -- (pq.north east);
\end{tikzpicture}

    \caption{Myšlenka důkazu korektnosti RSA přes moduly $p$, $q$ a~$pq$.}
    \label{fig:krypto-rsa-korektnost}
\end{figure}

\paragraph{Příklad s~malými čísly.}

\begin{example}
    Zvolme $p=5$, $q=11$ a~$s=27$. Potom
    \[
        n=55,
        \qquad
        \varphi(n)=(5-1)(11-1)=40.
    \]
    Protože $\operatorname{gcd}(27,40)=1$, můžeme nalézt $v$ splňující
    \[
        27v\bmod 40=1.
    \]
    Vyhovuje $v=3$, neboť $27\cdot 3=81\equiv 1\pmod{40}$. Veřejný klíč je $k_V=(3,55)$ a~soukromý klíč $k_S=(27,55)$.

    Alice zašifruje $x=7$:
    \[
        y=7^3\bmod 55=343\bmod 55=13.
    \]
    Bob poté získá
    \[
        13^{27}\bmod 55=7.
    \]
\end{example}

Malá čísla slouží pouze k~ručnímu ověření postupu. V~reálné kryptografii by tak malý modul záškodník rozložil okamžitě.

\subsection{Bezpečnost a~faktorizace}

Veřejný klíč obsahuje $n=pq$, ale nikoli samotná prvočísla $p$ a~$q$. Pokud záškodník rozklad zjistí, spočítá $\varphi(n)=(p-1)(q-1)$ a~z~veřejného exponentu $v$ odvodí soukromý exponent $s$. Efektivní faktorizace by tedy RSA prolomila.

Rozhodovací verzi problému můžeme zapsat takto:
\[
    \FACTOR:
    \quad
    \text{pro daná }n,b\in\N\text{ rozhodni, zda existuje }d
    \text{ takové, že }1<d<b\text{ a }d\mid n.
\]
Jak jsme viděli v~\smartref[sekci~][]{\showrefs}{sec:p_np}{dřívějším výkladu tříd P a~NP}, kladnou odpověď stačí doložit dělitelem $d$. Ověření nerovností i~rovnosti $n\bmod d=0$ proběhne v~polynomiálním čase, a~proto $\FACTOR\in\classNP$.

Naivní zkoušení dělitelů je vzhledem k~bitové délce $\sizeof{\binarycode{n}}$ exponenciální. Jsou známy podstatně rychlejší klasické algoritmy, žádný obecný polynomiální klasický algoritmus pro faktorizaci však znám není.

\importantbox{Znalost efektivní faktorizace stačí k~prolomení popsaného RSA. Není však dokázáno, že každý způsob, jak invertovat RSA, musí zároveň faktorizovat $n$. Bezpečnost RSA proto nelze jednoduše ztotožnit s~větou „faktorizace je těžká“.}

\warningbox{Uvedené „učebnicové RSA“ je deterministické a~bez doplňujícího kódování není bezpečným praktickým šifrovacím schématem. Skutečné systémy používají standardizované náhodné vycpávání a~oddělené postupy, například OAEP pro šifrování a~PSS pro podpis.}
